\begin{prob}
    In each of the problems below compute the expected time of the given code.
    State your answer using asymptotic notation. Show your work for this
    problem by stating what the different cases are, the probability of each
    case, and how long each case takes. Also show the calculation of the
    expected time.

    \begin{subprobset}
        \begin{subprob}
            \begin{minted}[autogobble]{python}
                def double_coin_toss(n):
                    coin_1 = np.random.rand() > 0.5
                    coin_2 = np.random.rand() > 0.5
                    if coin_1 and coin_2:
                        i = n
                        while i > 0:
                            print("We got two heads!")
                            i -= 1
                    else:
                        for i in range(n ** 2):
                            print("We didn't get two heads.")

            \end{minted}

            \begin{soln}
                $\Theta(n^{2})$.

                We can think of the two cases here as being 1) We got 2 heads, and 2) We did not get 2 heads.
                The probability of the first case is $1/4$ because the probability of getting each individual
                head is $1/2$ and we multiply $1/2 * 1/2$ for 2 coins, giving $1/4$.
                The probability of the second case is $3/4$ which comes from subtracting the probability of the
                first case from 1 giving $1 - 1/4 = 3/4$.

                In the first case, linear time is taken since we must loop over \python{n}.
                In the second case, quadratic time is taken since we must loop over \python{range(n**2)}.

                Therefore, the average time taken is:

                \begin{align*}
                    \frac{1}{4} \Theta(n) + \frac{3}{4} \Theta(n^2)\\
                    &=
                    \Theta(n) + \Theta(n^2)\\
                    &=
                    \Theta(n^{2})
                \end{align*}

            \end{soln}
        \end{subprob}

        \begin{subprob}
            \begin{minted}[autogobble]{python}
                def foo(n):
                    # randomly choose a number between 0 and n-1 in constant time
                    k = np.random.randint(n)

                    if k < np.sqrt(n):
                        for i in range(n**2):
                            print(i)
                    else:
                        print('Never mind...')

            \end{minted}

            \begin{soln}
                $\Theta(n\sqrt{n})$.

                We can think of the two cases here as being 1) $k$ is randomly chosen
                to be less than $\sqrt{n}$, and 2) $k$ is randomly chosen to be greater.
                The probability of the first case is $\sqrt{n} / n = 1/\sqrt{n}$, and
                the probability of the second case is $(n - \sqrt{n})/n = 1 - 1/\sqrt{n}$.

                In the first case, quadratic time is taken since we must loop over \python{range(n**2)}.
                In the second case, constant time is taken since we print \python{"Never mind..."}
                and exit.

                Therefore, the average time taken is:

                \begin{align*}
                    \frac{1}{\sqrt{n}} \Theta(n^2) + \left(1 - \frac{1}{\sqrt{n}}\right) \Theta(1)
                    =
                        \Theta(n\sqrt{n}) + \Theta(1)
                    =
                        \Theta(n\sqrt{n})
                \end{align*}

            \end{soln}
        \end{subprob}

        \begin{subprob}
            \begin{minted}[autogobble]{python}
                def baz(n):
                    # randomly choose a number between 0 and n-1 in constant time
                    x = np.random.randint(n)

                    for i in range(x):
                        for j in range(i):
                            print("Hello!")
            \end{minted}

            \textbf{Hint}: In class, we saw that the sum of the first $n$
            integers is $\frac{n(n+1)}{2}$. It turns out that the sum of the
            first $n$ integers squared (so, $1^2 + 2^2 + 3^2 + \ldots + n^2$)
            is $\frac{n(n+1)(2n+1)}{6}$. You can (and should) use this fact,
            but make sure to point it out when you do.


            \begin{soln}
                $\Theta(n^{2})$.

                Here, \python{x} is a random integer from 0 to $n$-1. We can
                think of each possible value of \python{x} as a different case.
                That is, in Case 0, we get 0 for \python{x}, in Case 1, we get
                1 for \python{x}, and so on up to Case $n-1$, where we get
                $n-1$ for \python{x}.

                Each case is equally likely, and there are $n$ of them, so the probability
                of each case is $1/n$.

                Next, we calculate the time taken in each case by counting the number
                of executions of the inner loop body. Consider Case $k$ (in which \python{x}
                turns out to be $k$). In this case, the nested loop becomes:

                \inputminted{python}{\thisdir/include-solution/loop.py}

                We have analyzed these sorts of dependent nested loops before
                in lecture, and we know that the number of executions of the
                inner loop body is equal to the sum $1 + 2 + 3 + \ldots + (k-1)$.
                This is an arithmetic sum, and it has a formula:
                $\frac{k(k-1)}{2}$. So, the number of executions in Case $k$ is
                $k(k-1)/2$, and therefore the time is proportional to this.

                Recalling that the expected time complexity has the formula:

                \begin{align*}
                    T_\text{expected}(n)
                    &=
                        \sum_{cases}
                        P(\text{case } k) \cdot T(\text{case } k)
                        \\
                \intertext{We fill in the probabilities and times:}
                    &= \sum_{k = 0}^{n-1} \frac{1}{n} \cdot \frac{k(k-1)}{2} \\
                    &= \frac{1}{2n} \sum_{k = 0}^{n-1} k(k-1) \\
                    &= \frac{1}{2n} \sum_{k = 0}^{n-1} k^2 - k \\
                    &= \frac{1}{2n} \left( \sum_{k = 0}^{n-1} k^2 - \sum_{k = 0}^{n-1} k \right) \\
                    \intertext{We recognize the second sum as an arithmetic sum, and we recognize the first sum from
                    the hint. Therefore:}
                    &= \frac{1}{2n} \left( \frac{n(n-1)(2n-1)}{6} - \frac{n(n-1)}{2} \right) \\
                    &= \frac{1}{2n} \left( \Theta(n^3) - \Theta(n^2) \right) \\
                    &= \Theta(n^2)
                \end{align*}

                The hardest part of this problem might have been figuring out $T(\text{case } k)$; that is,
                in counting the number of exeuptions of the inner loop body in Case $k$. If you got stuck here,
                or wanted to check your guess, remember that you can try it out in code:

                \inputminted{python}{\thisdir/include-solution/loop-count.py}

                This would print \python{36}, which matches the prediction of our
                formula when $k = 9$, since $9(9-1)/2 = 36$.

            \end{soln}

        \end{subprob}



    \end{subprobset}
\end{prob}
